%% ma: L11-B27-0018
%% lop: 11
%% chuong: 7
%% bai: 27
%% dang: 11.27.14
%% loai: DS
%% muc: [NB, VD, TH, TH]
%% dap_an: "ĐSĐS"
%% nguon: "(NBV)-ĐỀ SỐ 17-MỖI NGÀY 1 ĐỀ THI 2025.pdf (D17, MỖI NGÀY 1 ĐỀ - Đề số 17), Phần II Câu 1"
%% kien_thuc: "Bài 24 (góc giữa đường thẳng và mặt phẳng), Bài 26 (khoảng cách), Bài 27 (thể tích khối chóp)"
%% trang_thai: chinh-thuc
%% ghi_chu: "Đáp án gốc ĐSĐS đúng (giải độc lập: BC = √3, AC = 2, d(B,(SAC)) = √3/2, V = √3/3). Hình gốc không có. Đã duyệt 16/10/2026."
\begin{ex}
Cho hình chóp $S.ABC$ có đáy $ABC$ là tam giác vuông tại $B$ có $AB=1$, $\widehat{ACB}=30^\circ$. Biết $SA$ vuông góc với mặt đáy và $SA=2$. Gọi $H$ là hình chiếu của $A$ trên $SB$.
\choiceTF
{\True $d(A,SB)=AH$.}
{$d\bigl(B,(SAC)\bigr)=\dfrac{\sqrt3}{3}$.}
{\True $BC=\sqrt3$.}
{Thể tích khối chóp $S.ABC$ bằng $\dfrac{\sqrt3}{6}$.}
\loigiai{a) $H$ là hình chiếu của $A$ trên $SB$ nên $AH\perp SB$ và $d(A,SB)=AH$. Đúng.

b) $AC=\dfrac{AB}{\sin30^\circ}=2$. Gọi $K$ là hình chiếu của $B$ trên $AC$; do $SA\perp(ABC)$ nên $BK\perp SA$, suy ra $BK\perp(SAC)$ và $d\bigl(B,(SAC)\bigr)=BK=\dfrac{AB\cdot BC}{AC}=\dfrac{\sqrt3}{2}\ne\dfrac{\sqrt3}{3}$. Sai.

c) Tam giác $ABC$ vuông tại $B$ có $\widehat{ACB}=30^\circ$ nên $BC=\dfrac{AB}{\tan30^\circ}=\sqrt3$. Đúng.

d) $S_{ABC}=\dfrac12AB\cdot BC=\dfrac{\sqrt3}{2}$, $V=\dfrac13SA\cdot S_{ABC}=\dfrac13\cdot2\cdot\dfrac{\sqrt3}{2}=\dfrac{\sqrt3}{3}\ne\dfrac{\sqrt3}{6}$. Sai.}
\end{ex}
